\scanpage{46}
\deleted{Réciproquement, si ces conditions sont satisfaites, choisissant un gén $u$ de $Z$, on voit que} -- de sorte que $\{1,\sigma\}\mathbin{.}Z$ (produit $\frac12$ direct) opère sur $S$,\ldots \deleted{qui devient un polygône de}

On a évidemment

\uline{Prop 7} La donnée d'un \deleted{$n$-polygône} $\mathrm{SR}n\mathrm S$ dans $E$ équivaut à la donnée de a) b) ci-dessus (ss-grou \inserted{$Z$} cyclique d'ordre $n$ de $\mathrm{Aff}(E)$, et d'une orbite $S$ \deleted{de} $Z$ dans $E$). \deleted{faisant de \unsure{($E$ un type|$E$ un syst.)|les conditions}} \inserted{tels que les cond$^{\rm s}$} $1^\circ)$ et $2^\circ)$ ci-dessus soient vérifiées ($S$ non sur une droite, et condition de normalisation\ldots) -- et éventuellement la condition unimodulaire :

3) $Z$ formé d'endom unimodulaires.

\uline{Remarque} Suivant les cas ($p=\mathrm{car}\, k$, \unsure{etc|et|$n$)}, le groupe des autom d'un \deleted{syst} $\mathrm{SR}n\mathrm S$ est isomorphe à $\mathrm{Aff}^{\pm}(S)\simeq D_n$, ou à $\mathrm{Aff}(S)$. On peut se poser la question, dans l'un ou l'autre cas, de \deleted{savoir}
% archive sheet number: 45
